% Part: lambda-calculus
% Chapter: lambda-definability
% Section: partial-recursive-functions

\documentclass[../../../include/open-logic-section]{subfiles}

\begin{document}

\olfileid{lam}{ldf}{par}
\olsection{आंशिक पुनरावर्ती फलने \usetoken{S}{lambda definable} असतात}

मूलभूत फलनांपासून संयोजन, आदिम पुनरावर्तन आणि
अपरिबद्ध लघुतमीकरण वापरून मिळणारी फलने म्हणजे
आंशिक पुनरावर्ती फलने. सामान्य पुनरावर्ती फलनांपेक्षा
त्यांचा फरक असा की अपरिबद्ध शोधात वापरलेली फलने
नियमित असणे आवश्यक नाही. नियमिततेची अट नसल्याने
लघुतमीकरणाने परिभाषित केलेली फलने काही वेळा
अपरिभाषित राहू शकतात.

प्रथमदर्शनी, (सर्वत्र परिभाषित) सामान्य पुनरावर्ती
फलने !!{lambda definable} आहेत हे दाखवण्याच्या
पद्धतीनेच सर्व आंशिक पुनरावर्ती फलनेही
!!{lambda definable} आहेत हे सिद्ध करता येईल असे
वाटू शकते. उदाहरणार्थ, $f$ आणि $g$ ही अनुक्रमे
$F$ आणि~$G$ या पदांनी !!{lambda define}d
असतील, तर $f$ चे $g$ शी संयोजन
$\lambd[x][F (G x)]$ ने !!{lambda define}d होते.
पण फलने आंशिक असतील, तेव्हा अडचण येते.
$g(x)$ अपरिभाषित असल्यास $G x$ ला सामान्य रूप
नसते. बहुतेक वेळा मग $F (G x)$ लाही सामान्य रूप
नसते, आणि आपल्याला तेच हवे आहे. पण
$F$ हे $\lambd[x][\lambd[y][y]]$ असेल, तर
$F (G x)$ ला $\lambd[y][y]$ हे सामान्य रूप मिळते.

ही अडचण दूर करता येते. सर्व आंशिक पुनरावर्ती
फलने अशी !!{lambda define} करण्याचे मार्ग आहेत
की अपरिभाषित मूल्ये सामान्य रूप नसलेल्या पदांनी
दर्शवली जातात. मात्र हे मार्ग सामान्य पुनरावर्ती
फलनांसाठी वापरलेल्या पद्धतीपेक्षा काहीसे अधिक
गुंतागुंतीचे आणि कमी सहज आहेत. येथे पुरावा न देता
प्रमेय नोंदवतो:

\begin{thm}
  सर्व आंशिक पुनरावर्ती फलने !!{lambda definable} आहेत.
\end{thm}

\end{document}
